This chapter collects final thoughts. It is intended to provide an outlook for possible
further work. It also lists some points that went well, but also those where
potential for improvement was identified.

\subsection{Outlook}\label{sec:ausblick}

In principle, it can be confirmed that a \acrshort{tof} system, as imagined by the students,
can be developed and operated. In the following, a few possible improvements and further measurements are
listed.

\begin{enumerate}
      \item The optics should be better aligned. On the one hand, this would allow measurements to be taken over longer
            distances measurements over longer distances. Furthermore, it should also be possible to detect targets with
            non-reflective or weakly reflective surface properties.
      \item For measurements with the mirror, it would make sense to try a lens with a longer focal length (or the same
            lens at a slightly shorter distance to the \acrshort{pcb}). This would have the advantage that the mirror
            would not have to be aligned quite as precisely. The disadvantage would be a slightly lower photocurrent with
            parallel incident light.
      \item There is room for improvement in the control of the \acrshort{tdc}s. For example, a faster clocked
            \acrshort{mcu} could minimize the jitter at \lstinline|START| and \lstinline|STOP| signals, which would lead
            to measurement results with even less dispersion.
      \item Another way to minimize dispersion would be to use an external clock source.
      \item It may be possible to do without the artificially inserted delay of one clock cycle. This would be the case if
            there is already enough delay in the measurement path to exceed the minimum measurement duration of 12~ns.
            If not, the artificial delay could also be implemented by means of a separate circuit (with the lowest possible
            jitter) instead of the \acrshort{mcu} separate circuit (with the lowest possible jitter).
      \item The Nucleo used offers the option of configuring the rise time of the GPIO. With a faster rise time, it may be
            possible to improve the measurement results.
      \item Optical receiving path: Measurements should be carried out with different gains of the \acrshort{tia} to
            increase the distance with a mirror and to enable measurements against the wall or other targets.
            In this context, the switching threshold of the comparator should be tested more accurately and fine-tuned.
      \item Optical transmission path: It is worth experimenting with higher transmission power. Due to the 5~V supply
            of the laser diode, we are almost at the limit with the $5~\Omega$ series resistor. In a next \acrshort{pcb}
            version, it would be worthwhile to provide the option of a 12~V supply with a switch, analogous to how this is
            already implemented for the photodiode.
      \item Optical measurements: It is worth analyzing why the measurement results are distributed in two clusters
            (approx. 2 ns apart).
      \item For further measurements with a mirror, it is worth constructing a precise mechanical test setup.
            This would allow the mirror to be better fixed and adjusted. It is difficult to achieve reproducible results by hand.
\end{enumerate}

In conclusion, it is possible to perform a time-of-flight measurement with a TDC7200,
whereby the temporal resolution of the measurement itself is the limiting factor.

\subsection{Lessons Learned}

As the measurement results and the outlook already suggest, a little more time could have been invested primarily in
designing the optoelectronics. The initial rough estimates and calculations regarding the desired limitations of the
system should have been scrutinized more critically. Initially, the students' goal was to to take measurements at
distances of up to 10 cm with a resolution of about 10 cm. However, it turned out in the end that it was not the
temporal resolution but rather the distance and the strength of the measurement signal that was problematic.

The student team is satisfied with the fabrication of the \acrshort{tof} demonstrator. When using such small
SMD components, such as in the \acrshort{tia}, it is worth considering whether
machine assembly is preferable to manual assembly. Manual assembly probably saved hardly any
time, since it also took several hours, spread over several evenings.

\subsection{Acknowledgements}

We would like to take this opportunity to express our sincere thanks to the two lecturers, Prof. Guido Keel and Michael
Lehmann. During the project, they were always on hand to help us with advice and support and followed the project with
great interest. We would also like to thank them for the optical equipment loaned to us. This enabled us to achieve a
partial success in the optical measurement.

We would also like to thank Robin Burkard, who carried out a layout review for us. Thanks to his valuable tips, the
\acrshort{pcb} could be put into operation with as few glitches as possible.
